\chapter{Symmetric Krawtchouk Matrices}
\label{ch:16}

\section{Definition and small-order examples}
\label{sec:matrix-S-definition}
With the row/column convention of \eqref{eq:ency-Phi-generating}, put
\begin{equation}\label{eq:ency-S-definition}
S_N=\Phi_NB_N,\qquad
(S_N)_{nj}=\binom NjK_n(j;N).
\end{equation}
Thus $S_N$ is obtained by multiplying column $j$ of $\Phi_N$ by
$\binom Nj$, not by multiplying its row $j$.
Binomial duality gives
\[
\binom NjK_n(j;N)=\binom NnK_j(n;N),
\qquad S_N^T=S_N.
\]
For $N=2$,
\[
S_2=\begin{pmatrix}1&2&1\\2&0&-2\\1&-2&1\end{pmatrix}.
\]
For $N=4$, the column weights are $(1,4,6,4,1)$ and
\[
S_4=\begin{pmatrix}
1&4&6&4&1\\
4&8&0&-8&-4\\
6&0&-12&0&6\\
4&-8&0&8&-4\\
1&-4&6&-4&1
\end{pmatrix}.
\]
Although $\Phi_N^2=2^NI$, the inserted binomial weights prevent
identifying $S_N^2$ with $2^NI$. Nor is $S_N$ the same matrix as the
orthogonal normalization $U_N$ of
\eqref{eq:ency-U-normalization}. Complete tables of $A_N,R_N,L_N$,
$\Phi_N,S_N$ are in Appendix~\ref{app:matrix-tables}.

\section{The matrix $T$}
\label{sec:16-1}
% TODO: Develop this section; add verified references, examples, and exercises.

\noindent\textbf{THE MATRIX $T$ NEEDS TO BE DONE.}


\section{$S^2$ and trace formulas}
\label{sec:16-2}
Binomial duality yields the exact bilinear identity
\[
(S_N^2)_{jk}=\binom Nj\binom Nk\sum_{n=0}^NK_n(j;N)K_n(k;N).
\]
Parseval applied to the generating functions evaluates it as
\begin{equation}\label{eq:beta}
(S_N^2)_{jk}=\frac{4^N}{\pi}\binom Nj\binom Nk
\cos\frac{\pi(j-k)}2\,
\mathrm B\left(\frac{j+k+1}{2},\frac{2N-j-k+1}{2}\right).
\end{equation}
The beta integral follows by factoring $1\pm e^{i\theta}$ into half-angle sine and cosine factors. This identity is exact, not an asymptotic statement.

\section{Spectrum of $S$}
\label{sec:16-3}
% TODO: Develop this section; add verified references, examples, and exercises.

\noindent\textbf{SPECTRUM OF $S$ NEEDS TO BE DONE.}


\section{Symmetry sectors and block decompositions}
\label{sec:16-4}
Define $(Jf)_j=f_{N-j}$ and $(Ef)_j=(-1)^jf_j$. Then
\[
JS_N=S_NE,\qquad S_NJ=ES_N,
\]
so $S_N$ commutes with $J+E$. For the square, indices of different parity decouple. Conjugation by $\operatorname{diag}((-1)^{\lfloor j/2\rfloor})$ makes each parity block entrywise positive, by \eqref{eq:beta}.

\section{Beta-function kernel and spectral asymptotics}
\label{sec:16-5}
\begin{remark}[Working asymptotic program, not a completed theorem]
Put $\xi=(N-2j)/\sqrt N$ and $\eta=(N-2k)/\sqrt N$. Compact-window Stirling expansions of \eqref{eq:beta}, after sign removal and normalization by $N/8^N$, lead to the candidate kernel
\[
L(\xi,\eta)=\frac{1}{\sqrt2\pi^{3/2}}
\exp\left[-\frac38(\xi^2+\eta^2)+\frac14\xi\eta\right].
\]
Its Gaussian ground state is $\exp(-\xi^2/(2\sqrt2))$, with eigenvalue $\kappa=(\sqrt8-2)/\pi$. The proposed positive-eigenvalue asymptotic is
\[
\lambda_{\max}(S_N)\sim\sqrt{\sqrt8-2}\,
\frac{(\sqrt8)^N}{\sqrt{\pi N}}.
\]
Global operator convergence and control of the positive spectral branch still require proof. Fixed-degree Hermite limits alone do not justify this conclusion.
\end{remark}

